%TP genérico



%Configuracion del documento

%Tamaño de letra principal:

\documentclass[12pt]{article}


%Título y autor(es):

\title{Red de Hopfield}
\author{Sergio}


%Tamaño de página y los márgenes:
\usepackage[a4paper,headheight=16pt]{geometry}

\textwidth      =  450pt     %Ancho del cuerpo
\textheight     =  650pt     %Largo del cuerpo
\topmargin      =  0pt       %Agrega espacio en el margen superior
\oddsidemargin  =  0pt       %+ margen izquierdo en paginas impares
\evensidemargin =  0pt       %+ margen derecho en paginas impares


% castellano:
\usepackage[spanish]{babel}

%instalar texlive-lang-spanish
% Reconocer acentos y caracteres no ingleses:

\usepackage[utf8]{inputenc}


% Cabecera y pie de página personalizadas
\usepackage{fancyhdr}
\pagestyle{fancy}

% Hago que en la cabecera de página se muestre a la derecha la sección,
% y en el pie, en número de página a la derecha:

\renewcommand{\sectionmark}[1]{\markboth{}{\thesection\ \ #1}}
\lhead{}
\chead{}
\rhead{\rightmark}
\lfoot{}
\cfoot{}
\rfoot{\thepage}


%numeracion especial para tablas, figuras y ecuaciones
\usepackage{amsmath}
\numberwithin{equation}{section}
\numberwithin{figure}{section}
\numberwithin{table}{section}


%Agregar notas al pie en tablas:
\usepackage{threeparttable}


%Incluir Graficos
\usepackage{graphicx}


%Usar subfiguras: (al estilo Figura 2.3(b) )
\usepackage{subfigure}
\usepackage{float}

% Para esto es necesario texlive-latex-recommended o texlive-latex-extra

%Numero de figuras en negrita
\usepackage[hang,bf]{caption}


% Todas las imágenes están en el directorio tp-img:
\newcommand{\imgdir}{img}
\graphicspath{{\imgdir/}}

% uso de colores
\usepackage{color}
\definecolor{light-gray}{gray}{0.95}

%Para embeber código de lenguajes como Matlab, C, html, etc.

\usepackage{listings}
\lstset{ %
frameround=tttt,
language=bash,                  % choose the language of the code
basicstyle=\footnotesize,       % the size of the fonts that are used for the code
stepnumber=1,                   % the step between two line-numbers. If it is 1 each line will be numbered
backgroundcolor=\color{white},  % choose the background color. You must add \usepackage{color}
showspaces=false,               % show spaces adding particular underscores
showstringspaces=false,         % underline spaces within strings
showtabs=false,                 % show tabs within strings adding particular underscores
frame=single,           % adds a frame around the code
tabsize=4,          % sets default tabsize to 2 spaces
captionpos=b,           % sets the caption-position to bottom
breaklines=true,        % sets automatic line breaking
breakatwhitespace=false,    % sets if automatic breaks should only happen at whitespace
escapeinside={\%*}{*)}          % if you want to add a comment within your code
}

\lstdefinestyle{consola}
{basicstyle=\scriptsize\bf\ttfamily,
backgroundcolor=\color{light-gray},
}

%Para colocar hipervinculos

\usepackage[colorlinks=true,linkcolor=black,urlcolor=blue]{hyperref}

%Para hacer matrices y digramas de flujo sencillos

%\usepackage[all]{xy}
\usepackage{pdfpages}


%INICIO DEL DOCUMENTO

%-------------------------------------------------------------------------------



\begin{document}



% CARATULA

%-------------------------------------------------------------------------------



\begin{titlepage}

% Sin cabecera ni pie de página:

\thispagestyle{empty}



% Logo de la facu: 

	\includegraphics[width=5cm]{logo-facu}

\vfill



% Título:

\begin{center}
	\Huge{Descripción de objetos}\\

\end{center}

\vspace{4cm}



% Integrantes:

\large{

	\begin{tabbing}
		Sergio Hinojosa \hspace{1cm} \\84476 \\
	\end{tabbing}

}

\vfill

% Fecha o cuatrimestre:

%\flushright{2012}

\end{titlepage}




% Hago que las páginas se comiencen a contar a partir de aquí:

\setcounter{page}{1}



%INDICE

%-------------------------------------------------------------------------------

% Pongo el índice en una página aparte:

%\tableofcontents

%\newpage



%-------------------------------------------------------------------------------


% Inicio del TP:
\section{Instalación}

\subsection{En Ubuntu/Debian}

\subsection{NeoDB}
Instalar Neo:

\begin{lstlisting}[style=consola, language=bash]
$ apt-get install python-numpy python-pip
$ pip install quantities
$ apt-get install python-scipy
$ apt-get install libhdf5-serial-dev python-numexpr cython
$ pip install tables
$ pip install neo
\end{lstlisting}



\end{document}
